% 論文實驗章節示例（LaTeX 格式）
% 可以根據需要調整和整合到你的論文主體中

\section{Emotion Fusion Weight Optimization}
\label{sec:emotion_fusion_optimization}

\subsection{Problem Formulation}
\label{subsec:problem_formulation}

In our music recommendation system, accurate emotion detection from user posts is crucial. We propose a multi-signal fusion approach that combines three distinct emotion detection methods: (1) lexicon-based emotion detection using the NRC Emotion Lexicon, (2) emoji-based emotion detection, and (3) classifier-based emotion detection using a machine learning model. 

The key challenge is determining the optimal weight combination for fusing these three signals. Given that the emotion vector will be used for similarity computation with song emotion vectors, the weight selection significantly impacts recommendation quality.

\subsection{Fusion Strategy}
\label{subsec:fusion_strategy}

We employ a weighted linear combination approach:

\begin{equation}
E_{fused} = \alpha \cdot E_{lexicon} + \beta \cdot E_{emoji} + \gamma \cdot E_{classifier}
\end{equation}

where $\alpha + \beta + \gamma = 1.0$, and each emotion vector is an 8-dimensional probability distribution over emotion classes.

To address the issue of fallback signals (e.g., when no emojis are found, or when the lexicon returns only neutral), we implement a \textit{conditional fusion} mechanism that dynamically excludes invalid signals and renormalizes the weights:

\begin{algorithm}[h]
\caption{Conditional Emotion Fusion}
\begin{algorithmic}[1]
\Require Lexicon emotion $E_{lex}$, Emoji emotion $E_{emoji}$, Classifier emotion $E_{clf}$
\Require Weights $\alpha$, $\beta$, $\gamma$
\Ensure Fused emotion vector $E_{fused}$
\State Check if $E_{lex}$ is fallback (all neutral)
\State Check if $E_{emoji}$ is fallback (no emojis found)
\State Collect valid signals and their weights
\State Renormalize weights: $w'_i = w_i / \sum_j w_j$
\State Compute $E_{fused} = \sum_i w'_i \cdot E_i$
\State Normalize $E_{fused}$ to sum to 1.0
\end{algorithmic}
\end{algorithm}

\subsection{Experimental Setup}
\label{subsec:experimental_setup}

\subsubsection{Dataset}
We evaluate on a labeled dataset of 10,000 social media posts, evenly distributed across 8 emotion categories (approximately 12-13\% per class): joy, anger, fear, sadness, surprise, disgust, excitement, and neutral.

\subsubsection{Weight Combinations}
We systematically test 50+ weight combinations, including:
\begin{itemize}
    \item Pure method baselines: $(0, 0, 1.0)$, $(1.0, 0, 0)$, $(0, 1.0, 0)$
    \item Classifier-dominant combinations: $\gamma \geq 0.85$, with fine-grained steps of 0.01-0.02
    \item Balanced combinations with similar weights
    \item Method-specific priority combinations
\end{itemize}

\subsubsection{Evaluation Metrics}
Given that emotion vectors are used for similarity computation with song emotion vectors, we employ a comprehensive set of metrics:

\textbf{Classification Metrics:}
\begin{itemize}
    \item Accuracy: proportion of correctly predicted labels
    \item F1-score (macro and weighted): class-level and sample-weighted averages
\end{itemize}

\textbf{Vector Similarity Metrics:}
\begin{itemize}
    \item \textbf{Top-2 Accuracy}: whether the true label is in the top-2 predicted emotions (primary metric for similarity-based matching)
    \item Mean Cosine Similarity: cosine similarity between predicted and true one-hot vectors
    \item Mean Cross-Entropy: cross-entropy loss
    \item Mean Squared Error (MSE): L2 distance
    \item KL Divergence: distribution distance
\end{itemize}

\textbf{Distribution Analysis:}
\begin{itemize}
    \item Neutral overprediction rate: difference between predicted and true neutral proportions
\end{itemize}

\subsection{Results}
\label{subsec:results}

\subsubsection{Baseline Methods}

Table~\ref{tab:baseline_results} shows the performance of individual methods:

\begin{table}[h]
\centering
\caption{Baseline Method Performance}
\label{tab:baseline_results}
\begin{tabular}{lcccc}
\toprule
Method & Accuracy & Top-2 Acc. & F1 (macro) & Cosine Sim. \\
\midrule
Pure Classifier & \textbf{0.7439} & \textbf{0.8604} & \textbf{0.7288} & \textbf{0.3737} \\
Pure Lexicon & 0.3144 & 0.5135 & 0.2697 & 0.2476 \\
Pure Emoji & 0.2916 & 0.4294 & 0.2633 & 0.2355 \\
\bottomrule
\end{tabular}
\end{table}

The classifier significantly outperforms lexicon and emoji methods, suggesting it should play a dominant role in the fusion.

\subsubsection{Top Performing Combinations}

Table~\ref{tab:top_combinations} presents the top-5 weight combinations ranked by Top-2 Accuracy, which is our primary metric for similarity-based matching:

\begin{table}[h]
\centering
\caption{Top-5 Weight Combinations (Ranked by Top-2 Accuracy)}
\label{tab:top_combinations}
\begin{tabular}{lcccccc}
\toprule
Rank & Weights & Top-2 & Accuracy & F1 (macro) & Neutral \\
    & (Lex, Emoji, Clf) & Acc. & & & Overpred. \\
\midrule
1 & (0.03, 0.07, 0.90) & \textbf{0.8650} & 0.7559 & 0.7396 & +3.09\% \\
2 & (0.03, 0.04, 0.93) & 0.8649 & 0.7532 & 0.7376 & +3.19\% \\
3 & (0.03, 0.05, 0.92) & 0.8645 & 0.7546 & 0.7388 & +3.13\% \\
4 & (0.04, 0.06, 0.90) & 0.8644 & \textbf{0.7565} & \textbf{0.7403} & \textbf{+2.99\%} \\
5 & (0.01, 0.04, 0.95) & 0.8639 & 0.7506 & 0.7350 & +3.29\% \\
\midrule
Pure Clf. & (0, 0, 1.0) & 0.8604 & 0.7439 & 0.7288 & +3.04\% \\
\bottomrule
\end{tabular}
\end{table}

\subsection{Analysis and Discussion}
\label{subsec:analysis}

\subsubsection{Why Classifier Dominates}

The classifier achieves superior performance due to:
\begin{enumerate}
    \item \textbf{Training data}: Trained on 11,384 labeled posts, achieving 84.28\% test accuracy
    \item \textbf{Contextual understanding}: Captures semantic, syntactic, and contextual patterns that lexicon and emoji methods miss
    \item \textbf{Generalization}: Better handles new words, internet slang, and implicit expressions
\end{enumerate}

\subsubsection{Why Lexicon and Emoji Have Limited Contribution}

\textbf{Lexicon Limitations:}
\begin{itemize}
    \item \textbf{Coverage}: Only matches words in the dictionary, missing new terms and slang
    \item \textbf{Context}: Cannot handle negations, sarcasm, or complex expressions
    \item \textbf{Performance}: Pure lexicon accuracy is only 31.44\%
\end{itemize}

\textbf{Emoji Limitations:}
\begin{itemize}
    \item \textbf{Coverage}: Only 68.64\% of posts contain covered emojis
    \item \textbf{Error propagation}: Long conversion chain (emoji → phrase → NRC) accumulates errors
    \item \textbf{Fallback}: Many emojis lack phrase mappings, defaulting to neutral
    \item \textbf{Performance}: Pure emoji accuracy is only 29.16\%, with +29.09\% neutral overprediction
\end{itemize}

\subsubsection{Conditional Fusion Impact}

The conditional fusion mechanism significantly improves results by:
\begin{itemize}
    \item Excluding fallback signals (lexicon-only neutral, emoji-not-found neutral)
    \item Preventing low-quality signals from diluting the fusion
    \item Dynamically renormalizing weights based on available signals
\end{itemize}

Without conditional fusion, emoji fallback cases would contribute neutral vectors with their full weight, degrading overall performance.

\subsubsection{Weight Selection Insights}

Our experiments reveal:
\begin{enumerate}
    \item \textbf{High-quality signals dominate}: The classifier's high accuracy warrants the highest weight ($\gamma \geq 0.85$)
    \item \textbf{Low-quality signals used sparingly}: Lexicon and emoji weights should be small ($\alpha + \beta \leq 0.15$)
    \item \textbf{Diminishing returns}: Increasing $\alpha$ or $\beta$ beyond certain thresholds provides minimal or negative gains
\end{enumerate}

\subsection{Selected Configuration}
\label{subsec:selected_config}

Based on comprehensive evaluation, we select:

$$\alpha = 0.04, \quad \beta = 0.06, \quad \gamma = 0.90$$

This combination achieves:
\begin{itemize}
    \item Top-2 Accuracy: 0.8644 (only 0.0006 below the maximum)
    \item Accuracy: 0.7565 (highest among top combinations)
    \item F1 (macro): 0.7403 (highest)
    \item Neutral overprediction: +2.99\% (lowest, most balanced)
\end{itemize}

This configuration balances multiple metrics while prioritizing Top-2 Accuracy, which is crucial for similarity-based song matching in our recommendation system.

